\subsection{The Galois group}

DEFINITION 2. --- Let N be a Galois extension of the field K. The group of all K-automorphisms of N will be called the Galois group of N over K and denoted by Gal(N/K).

Let N be a finite Galois extension of K. Then N is a finite separable and quasi-Galois extension of K. Therefore (V, p.~32, Prop. 4 and V, p.~54, Prop. 3), the order of Gal(N/K) is equal to {[}N : K{]}. We shall prove later that if N is a Galois extension of K such that Gal(N/K) is finite, then N is of finite degree over K (V, p.~66, Th. 3).

Let \(\Omega\) be an algebraically closed extension of K and let A be the set of roots in \(\Omega\) of a separable polynomial \(f \in \mathbf{K}[\mathbf{X}]\) . Then the field \(\mathbf{N} = \mathbf{K}(\mathbf{A})\) is a Galois extension of K. It is clear that every K-automorphism of N leaves A stable, and since A generates N over K, the mapping \(\sigma \mapsto \sigma | \mathbf{A}\) is an isomorphism of \(\operatorname{Gal}(\mathbf{N} / \mathbf{K})\) onto a subgroup \(\Gamma\) of the symmetric group \(\mathfrak{S}_{\mathbf{A}}\) of the set A, which will be called the Galois group of the polynomial \(f\) . From Remark 3 (V, p.~55) it follows that if \(x\) and \(y\) belong to A, the following properties are equivalent:

\begin{enumerate}
\def\labelenumi{\alph{enumi})}
\item
  \(x\) and \(y\) are conjugate over \(\mathbf{K}\) ,
\item
  \(x\) and \(y\) belong to the same orbit under \(\Gamma\) ,
\item
  \(x\) and \(y\) are roots of the same irreducible factor of \(f\) .
\end{enumerate}

In particular \(f\) is irreducible if and only if \(\mathbf{A}\) is non-empty and \(\Gamma\) operates transitively on \(\mathbf{A}\) .

Examples. --- 1) Suppose that the characteristic of K is different from 2 and let N be a quadratic extension of K. If \(x \in \mathbf{N} - \mathbf{K}\) , then we have \(\mathbf{N} = \mathbf{K}(x)\) and the minimal polynomial of \(x\) over K is of the form \(f(\mathbf{X}) = \mathbf{X}^2 - a\mathbf{X} + b\) , with \(a, b \in \mathbf{K}\) . We thus have \(f(\mathbf{X}) = (\mathbf{X} - x)(\mathbf{X} - y)\) where \(y = a - x\) , so \(y\) is conjugate to \(x\) ; since \(f(\mathbf{X})\) is separable, the extension N is Galois. The group \(\operatorname{Gal}(\mathbf{N}/\mathbf{K})\) has two elements which induce the two permutations of the set \(\{x, y\}\) .

2) Let \(f = \mathbf{X}^3 + \mathbf{X}^2 - 2\mathbf{X} - 1 \in \mathbf{Q}[\mathbf{X}]\) . The polynomial \(f\) is irreducible, for if not, then it would have a root \(x \in \mathbf{Q}\) ; on writing \(x = a / b\) with \(a, b \in \mathbf{Z}\) , \(a\) and \(b\) relatively prime, we should find \(a(a^2 + ab - 2b^2) = b^3\) and \(a^3 = b(b^2 + 2ab - a^2)\) ; but this implies that \(a\) divides \(b\) and \(b\) divides \(a\) , hence \(x = \pm 1\) , which is impossible. Put \(\xi = e^{2\pi i / 7} \in \mathbf{C}\) ; then the polynomial \(f\) has the roots \(\alpha = \xi + \xi^{-1}, \beta = \xi^2 + \xi^{-2}, \gamma = \xi^3 + \xi^{-3}\) . We have \(\beta = \alpha^2 - 2\) and \(\gamma = \alpha^3 - 3\alpha\) , hence the extension \(\mathbf{Q}(\alpha)\) is Galois over \(\mathbf{Q}\) . The Galois group of \(\mathbf{Q}(\alpha)\) over \(\mathbf{Q}\) is cyclic of order 3 and is generated by an element \(\sigma\) such that \(\sigma(\alpha) = \beta, \sigma(\beta) = \gamma, \sigma(\gamma) = \alpha.\)

\begin{enumerate}
\def\labelenumi{\arabic{enumi})}
\setcounter{enumi}{2}
\item
  Suppose that K = Q and take \(f = X^{3} - 2\) . Using the decomposition of integers into products of prime factors (I, p.~51), we easily see that 2 is not the cube of an element of Q. Hence the polynomial f is irreducible, for if not, it would have a root in Q. Let \(A = \{x_{1}, x_{2}, x_{3}\}\) be the set of roots of f in \(\Omega\) and \(\Gamma\) the Galois group of f.~It acts transitively on A; its order is therefore divisible by three. On the other hand, the quotient \(j = \frac{x_{2}}{x_{1}}\) is different from 1 and we have \(j^{3} = 1\) . Hence j satisfies the relation \(j^{2} + j + 1 = 0\) ; now the polynomial \(T^{2} + T + 1 = \left(T + \frac{1}{2}\right)^{2} + \frac{3}{4}\) has no roots in Q, which shows (Example 1) that \([Q(j): Q] = 2\) . Thus \([N: Q]\) is divisible by 2, and it follows that the order of \(\Gamma\) is divisible by 6. Since \(\Gamma\) is contained in the group \(S_{A}\) of order 6, we have \(\Gamma = S_{A}\) .
\item
  Suppose that K is of characteristic \(p \neq 0\) and let \(\mathrm{K}(\mathrm{T})\) be the field of rational fractions and \(U = T^{p} - T\) . Put \(E = K[U]\) and \(F = K[T]\) , then the polynomial \(f(X) = X^{p} - X - U\) of \(E[X]\) has the roots \(T, T + 1, \ldots, T + p - 1\) in F. Let \(\sigma\) be the K-automorphism of F such that \(\sigma(T) = T + 1\) . We have \(\sigma^{i}(T) = T + i\) and \(\sigma(U) = U\) . The group \(G = \{1, \sigma, \ldots, \sigma^{p-1}\}\) is cyclic of order p, and its field of invariants contains E; since \([F: E] \leqslant p\) , Dedekind's theorem (V, p.~27, Cor. 2) implies that E is the field of invariants of G and \([F: E] = p\) . The polynomial f is thus irreducible in \(E[X]\) ; the extension F of E is Galois, its Galois group G is cyclic of order p, and the group \(\Gamma\) is the group of cyclic permutations of T, \(T + 1, \ldots, T + p - 1\) .
\end{enumerate}

For a generalization of this example see V, p.~93, Example 2.

\begin{enumerate}
\def\labelenumi{\arabic{enumi})}
\setcounter{enumi}{4}
\tightlist
\item
  Let \(\mathbf{F} = \mathbf{K}(\mathbf{X}_1, \ldots, \mathbf{X}_n)\) be the field of rational fractions in \(n\) indeterminates \(\mathbf{X}_1, \ldots, \mathbf{X}_n\) with coefficients in \(\mathbf{K}\) . Put
\end{enumerate}

\[
s _ {k} = \sum_ {1 \leqslant i _ {1} <   \dots <   i _ {k} \leqslant n} X _ {i _ {1}} \dots X _ {i _ {k}}
\]

for \(1 \leqslant k \leqslant n\) and \(\mathrm{E} = \mathrm{K}(s_1, \ldots, s_n)\) ; we denote by \(f(\mathrm{T})\) the polynomial

\[
\mathrm{T} ^ {n} - s _ {1} \mathrm{T} ^ {n - 1} + \dots + (- 1) ^ {n} s _ {n}.
\]

We have \(f(\mathrm{T}) = \prod_{i=1}^{n} (\mathrm{T} - \mathrm{X}_i)\) , so that \(\mathbf{F}\) is a splitting field of the separable polynomial \(f(\mathrm{T}) \in \mathrm{E}[\mathrm{T}]\) . Further, for every permutation \(\sigma \in \mathfrak{S}_n\) there exists a unique \(\mathbf{K}\) -automorphism \(h_\sigma\) of \(\mathbf{F}\) such that \(h_\sigma(\mathrm{X}_i) = \mathrm{X}_{\sigma(i)}\) for \(1 \leqslant i \leqslant n\) ; we have \(h_\sigma(s_k) = s_k\) for \(1 \leqslant k \leqslant n\) , hence \(h_\sigma\) is an E-automorphism of \(\mathbf{F}\) . In other words, \(\mathbf{F}\) is a Galois extension of \(\mathbf{E}\) and the restriction to the set of roots \(\{\mathbf{X}_1, ..., \mathbf{X}_n\}\) of \(f(\mathrm{T})\) defines an isomorphism of \(\operatorname{Gal}(\mathrm{F}/\mathrm{E})\) onto the group \(\mathfrak{S}_n\) . In particular, \(\mathbf{E}\) consists of rational fractions \(f\) such that

\[
f \left(\mathrm{X} _ {\sigma (1)}, \dots , \mathrm{X} _ {\sigma (n)}\right) = f \left(\mathrm{X} _ {1}, \dots , \mathrm{X} _ {n}\right)
\]

for every \(\sigma \in \mathfrak{S}_n\) (cf.~IV, p.~67, Cor.).

\begin{enumerate}
\def\labelenumi{\arabic{enumi})}
\setcounter{enumi}{5}
\tightlist
\item
  Suppose that \(f\) is monic of degree \(> 0\) and \(K\) is of characteristic \(\neq 2\) . Define on \(A\) a total ordering, denoted by \(\leqslant\) , and put \(\delta(f) = \prod_{\alpha < \beta} (\beta - \alpha)\) ,
\end{enumerate}

\((\alpha, \beta) \in \mathbf{A} \times \mathbf{A}\) , and for each \(\sigma \in \mathfrak{S}_{\mathbf{A}}\) put \(\delta_{\sigma}(f) = \prod_{\alpha < \beta} (\sigma(\beta) - \sigma(\alpha))\) . We have

\(\delta_{\sigma}(f) = \varepsilon(\sigma)\ \delta(f)\) , where \(\varepsilon(\sigma)\) is the signature of \(\sigma\) (I, p.~64) and \(\delta(f) \neq 0\) . For every \(\tau \in \operatorname{Gal}(\mathbf{N}/\mathbf{K})\) we have \(\tau(\delta(f)) = \delta_{\tau|_{\mathbf{A}}}(f)\) . Hence \(\Gamma\) is contained in the alternating group \(\mathfrak{A}_{\mathbf{A}}\) if and only if \(\delta(f) \in \mathbf{K}\) . Moreover \(\delta(f)^2 = \prod_{\alpha < \beta} (\beta - \alpha)^2 =\)

\(d(f)\) is the discriminant of the polynomial \(f\) (IV, p.~81) . Hence \(\Gamma \subset \mathfrak{A}_{\mathbf{A}}\) if and only if \(d(f)\) is the square of an element of \(\mathbf{K}\) . Thus in Example 2 we have \(d(f) = 49 = 7^2\) and in Example 3, \(d(f) = -108\) (IV, p.~85).

Let N be a Galois extension of K and let L be a subextension of N which is Galois over K. Every K-automorphism \(\sigma\) of N induces a K-automorphism \(\sigma_{L}\) of L (V, p.~55, Remark 1). Therefore the mapping \(\sigma \mapsto \sigma_{L}\) is a homomorphism of \(\operatorname{Gal}(N/K)\) into \(\operatorname{Gal}(L/K)\) , called the restriction homomorphism.

PROPOSITION 3. --- The restriction homomorphism of Gal(N/K) into Gal(L/K) is surjective.

More generally, consider two subextensions L and L' of N, and a K-isomorphism u of L onto L'. Choose an algebraic closure \(\Omega\) of K containing N as subextension (V, p.~23, Th. 2). There exists a K-automorphism v of \(\Omega\) which coincides with u on L (V, p.~52, Prop. 1), and since N is a quasi-Galois extension of K, v induces a K-automorphism \(\sigma\) of N (V, p.~55, Remark 1). In other words, the element \(\sigma\) of Gal(N/K) coincides with u on L.

